\documentclass{beamer}
\usepackage{amsfonts,amsmath,oldgerm}
\usetheme{sintef} % options: nocorners, nopicturecorners, noslidenumber

\usetikzlibrary{positioning,arrows.meta,fit,backgrounds} % for the diagrams section
% "visible on=<2->": the element takes its space on every slide but is shown
% only from slide 2, so the diagram does not shift while it is revealed.
% (for pictures inside a node, use {\visible<2->{\includegraphics...}} instead,
%  with the extra braces)
\tikzset{
  invisible/.style={opacity=0, text opacity=0},
  visible on/.style={alt={#1{}{invisible}}},
  alt/.code args={<#1>#2#3}{\alt<#1>{\pgfkeysalso{#2}}{\pgfkeysalso{#3}}},
}


\newcommand{\testcolor}[1]{\colorbox{#1}{\textcolor{#1}{test}}~\texttt{#1}}

\usefonttheme[onlymath]{serif}

% \titlebackground*{assets/background} % optional: replaces the faded emblem

\newcommand{\hrefcol}[2]{\textcolor{cyan}{\href{#1}{#2}}}

\title{Sapienza Beamer Presentation Theme}
\subtitle{Using \LaTeX\ to prepare slides}
\course{PhD in Natural Language Processing}
\author{Sapienza NLP Lab}
\date{2026-11-10}
\footertext{Sapienza Beamer Presentation Theme} % optional: omit for no footer text
% \cornerscale{1} % optional: size of the corner triangles (default 0.7)

\begin{document}
\maketitle

\begin{frame}

This template is a fork of \hrefcol{https://github.com/andrea-gasparini/sapienza-beamer-template}{sapienza-beamer-template} by \hrefcol{mailto:andrea@gasparini.cloud}{Andrea Gasparini}, restyled after the Sapienza NLP slide layout.

\vspace{\baselineskip}

The original template is based on \hrefcol{https://www.overleaf.com/latex/templates/sintef-presentation/jhbhdffczpnx}{SINTEF Presentation} from \hrefcol{mailto:federico.zenith@sintef.no}{Federico Zenith} and its derivation \hrefcol{https://github.com/TOB-KNPOB/Beamer-LaTeX-Themes}{Beamer-LaTeX-Themes} from Liu Qilong

\vspace{\baselineskip}

In the following you find a brief introduction on how to use \LaTeX\ and the beamer package to prepare slides, based on the one written by \hrefcol{mailto:federico.zenith@sintef.no}{Federico Zenith} for \hrefcol{https://www.overleaf.com/latex/templates/sintef-presentation/jhbhdffczpnx}{SINTEF Presentation}

% This template is released under \hrefcol{https://creativecommons.org/licenses/by-nc/4.0/legalcode}{Creative Commons CC BY 4.0} license

\end{frame}

\section{Introduction}

\begin{frame}{Beamer for SINTEF slides}
\begin{itemize}
\item We assume you can use \LaTeX; if you cannot,
\hrefcol{http://en.wikibooks.org/wiki/LaTeX/}{you can learn it here}
\item Beamer is one of the most popular and powerful document
classes for presentations in \LaTeX
\item Beamer has also a detailed
\hrefcol{http://www.ctan.org/tex-archive/macros/latex/contrib/beamer/doc/beameruserguide.pdf}{user
 manual}
\item Here we will present only the most basic features to get you up to speed
\end{itemize}
\end{frame}

\begin{frame}{Beamer vs. PowerPoint}
Compared to PowerPoint, using \LaTeX\ is better because:
\begin{itemize}
\item It is not What-You-See-Is-What-You-Get, but
What-You-\emph{Mean}-Is-What-You-Get:\\
you write the content, the computer does the typesetting
\item Produces a \texttt{pdf}: no problems with fonts, formulas,
      program versions
\item Easier to keep consistent style, fonts, highlighting, etc.
\item Math typesetting in \TeX\ is the best:
\begin{equation*}
\mathrm{i}\,\hslash\frac{\partial}{\partial t} \Psi(\mathbf{r},t) =
-\frac{\hslash^2}{2\,m}\nabla^2\Psi(\mathbf{r},t)
+ V(\mathbf{r})\Psi(\mathbf{r},t)
\end{equation*}

\end{itemize}
\end{frame}

\begin{frame}[fragile]{Getting Started}
\framesubtitle{Selecting the SINTEF Theme}
To start working with \texttt{sintefbeamer}, start a \LaTeX\ document with the
preamble:
\begin{block}{Minimum SINTEF Beamer Document}
\verb|\documentclass{beamer}|\\
\verb|\usetheme{sintef}|\\
\verb|\begin{document}|\\
\verb|\begin{frame}{Hello, world!}|\\
\verb|\end{frame}|\\
\verb|\end{document}|\\
\end{block}
\end{frame}

\begin{frame}[fragile]{Title page}
To set a typical title page, you call some commands in the preamble:
\begin{block}{The Commands for the Title Page}
\begin{verbatim}
\title{Sample Title}
\subtitle{Sample subtitle}
\author{First Author, Second Author}
\date{\today} % Can also be (ab)used for conference name &c.
\end{verbatim}
\end{block}
You can then write out the title page with \verb|\maketitle|.

To set a \textbf{background image} use the \verb|\titlebackground| command 
before \verb|\maketitle|; its only argument is the name (or path) of a graphic 
file.

If you use the \textbf{starred version} \verb|\titlebackground*|, the image 
will be clipped to a split view on the right side of the title slide.

\end{frame}

\begin{frame}[fragile]{Writing a Simple Slide}
\framesubtitle{It's really easy!}
\begin{itemize}[<+->]
\item A typical slide has bulleted lists
\item These can be uncovered in sequence
\end{itemize}
\begin{block}{Code for a Page with an Itemised List}<+->
\begin{verbatim}
\begin{frame}{Writing a Simple Slide}
  \framesubtitle{It's really easy!}
  \begin{itemize}[<+->]
    \item A typical slide has bulleted lists
    \item These can be uncovered in sequence
  \end{itemize}\end{frame}
\end{verbatim}
\end{block}
\end{frame}

\section{Personalization}

\footlinecolor{sintefyellow}
\begin{frame}[fragile]{Changing Slide Style}
\begin{itemize}
\item You can select the white or \textit{maincolor} \textbf{slide style} \emph{in the 
preamble} with \verb|\themecolor{white}| (default) or \verb|\themecolor{main}|
      \begin{itemize}
      \item You should \emph{not} change these within the document: Beamer does 
      not like it
      \item If you \emph{really} must, you may have to add 
      \verb|\usebeamercolor[fg]{normal text}| in the slide
      \end{itemize}
\item You can change the \textbf{footline colour} with 
\verb|\footlinecolor{color}|
      \begin{itemize}
      \item Place the command \emph{before} a new \verb|frame|
      \item There are four ``official'' colors: 
      \testcolor{maincolor}, \testcolor{sintefyellow}, 
      \testcolor{sintefgreen}, \testcolor{sintefdarkgreen}
      \item Default is no footline; you can restore it with 
      \verb|\footlinecolor{}|
      \item Others may work, but no guarantees!
      \item Should \emph{not} be used with the \verb|maincolor| theme!
      \end{itemize}
\end{itemize}
\end{frame}

\begin{frame}[fragile]{Blocks}
\begin{columns}
\begin{column}{0.4\textwidth}
\begin{block}{Standard Blocks}
These have a color coordinated with the footline (and grey in the blue theme)
\begin{verbatim}
\begin{block}{title}
content...
\end{block}
\end{verbatim}
\end{block}
\end{column}
\begin{column}{0.5\textwidth}
\begin{colorblock}[black]{sinteflightgreen}{Colour Blocks}
Similar to the ones on the left, but you pick the colour. Text will be white by 
default, but you may set it with an optional argument.
\small
\begin{verbatim}
\begin{colorblock}[black]
{sinteflightgreen}{title}
    content...
\end{colorblock}
\end{verbatim}
\end{colorblock}
The ``official'' colours of colour blocks are: \testcolor{sinteflilla}, 
\testcolor{maincolor}, \testcolor{sintefdarkgreen}, and 
\testcolor{sintefyellow}.
\end{column}
\end{columns}
\end{frame}

\footlinecolor{}
\begin{frame}[fragile]{Using Colours}
\begin{itemize}[<alert@2>]
  \item You can use colours with the
        \verb|\textcolor{<color name>}{text}| command
  \item The colours are defined in the \texttt{sintefcolor} package:
  \begin{itemize}
  \item Primary colours: \testcolor{maincolor} and its sidekick 
  \testcolor{sintefgrey}
  \item Three shades of green: \testcolor{sinteflightgreen}, 
  \testcolor{sintefgreen}, \testcolor{sintefdarkgreen}
  \item Additional colours: \testcolor{sintefyellow}, \testcolor{sintefred}, 
        \testcolor{sinteflilla}
        \begin{itemize}
        \item These may be shaded---see the \verb|sintefcolor| documentation or 
        the \hrefcol{https://sintef.sharepoint.com/sites/stottetjenester/%
        kommunikasjon/grafisk-profil-new/Sider/default.aspx}{SINTEF profile 
        manual}
        \end{itemize}
  \end{itemize}
  \item Do \emph{not} abuse colours: \verb|\emph{}| is usually enough
  \item Use \verb|\alert{}| to bring the \alert<2->{focus} somewhere
  \item<2- | alert@2> If you highlight too much, you don't highlight at all!
\end{itemize}
\end{frame}

\begin{frame}[fragile]{Adding images}
\begin{columns}
\begin{column}{0.6\textwidth}
Adding images works like in normal \LaTeX:
\begin{block}{Code for Adding Images}
\begin{verbatim}
\usepackage{graphicx}
    % ...
\includegraphics[width=\textwidth]
{assets/logo_RGB}
\end{verbatim}
\end{block}
\end{column}
\begin{column}{0.3\textwidth}
\includegraphics[width=\textwidth]
{assets/logo_RGB}
\end{column}
\end{columns}
\end{frame}

\begin{frame}[fragile]{Splitting in Columns}
Splitting the page is easy and common;
typically, one side has a picture and the other text:
\begin{columns}
\begin{column}{0.6\textwidth}
This is the first column
\end{column}
\begin{column}{0.3\textwidth}
And this the second
\end{column}
\end{columns}
\begin{block}{Column Code}
\begin{verbatim}
\begin{columns}
    \begin{column}{0.6\textwidth}
        This is the first column
    \end{column}
    \begin{column}{0.3\textwidth}
        And this the second
    \end{column}
    % There could be more!
\end{columns}
\end{verbatim}
\end{block}
\end{frame}

\begin{chapter}[assets/background_alternative]{}{Special Slides}
\begin{itemize}
\item Chapter slides
\item Side-picture slides
\end{itemize}
\end{chapter}

\footlinecolor{sintefred}
\begin{frame}[fragile]{Chapter slides}
\begin{itemize}
\item Similar to \verb|frame|s, but with a few more options
\item Opened with \verb|\begin{chapter}[<image>]{<color>}{<title>}|
\item Image is optional, colour and title are mandatory
\item There are seven ``official'' colours: \testcolor{maincolor}, 
\testcolor{sintefdarkgreen}, \testcolor{sintefgreen}, 
\testcolor{sinteflightgreen}, \testcolor{sintefred}, \testcolor{sintefyellow}, 
\testcolor{sinteflilla}.
      \begin{itemize}
      \item Strangely enough, these are \emph{more} than the official colours 
      for the footline.
      \item It may still be a nice touch to change the footline of following 
      slides to the same color of a chapter slide. Your choice.
      \end{itemize}
\item Otherwise, \verb|chapter| behaves just like \verb|frame|.
\end{itemize}
\end{frame}

\begin{sidepic}{assets/background_alternative}{Side-Picture Slides}
\begin{itemize}
\item Opened with \texttt{$\backslash$begin\{sidepic\}\{<image>\}\{<title>\}}
\item Otherwise, \texttt{sidepic} works just like \texttt{frame}
\end{itemize}
\end{sidepic}

\footlinecolor{}
\begin{frame}
\frametitle{Fonts}
\begin{itemize}
\item The paramount task of fonts is being readable
\item There are good ones...
  \begin{itemize}
  \item {\textrm{Use serif fonts only with high-definition projectors}}
  \item {\textsf{Use sans-serif fonts otherwise (or if you simply prefer 
them)}}
  \end{itemize}
\item ... and not so good ones:
  \begin{itemize}
  \item {\texttt{Never use monospace for normal text}}
  \item {\frakfamily Gothic, calligraphic or weird fonts: should always: be
  avoided}
\end{itemize}
\end{itemize}
\end{frame}

\begin{frame}[fragile]{Look}
\begin{itemize}
\item To insert a final slide with the title and final thanks, use \verb|\backmatter|.
      \begin{itemize}
      \item The title also appears in footlines along with the author name, you can change this text with \verb|\footlinepayoff|
      \item You can remove the title from the final slide with \verb|\backmatter[notitle]|
      \end{itemize}
\item The aspect ratio defaults to 16:9, and you should not change it to 4:3
      for old projectors as it is inherently impossible to perfectly convert a 
      16:9 presentation to 4:3 one; spacings \emph{will} break
      \begin{itemize}
      \item The \texttt{aspectratio} argument to the \texttt{beamer} class is
            overridden by the SINTEF theme
      \item If you \emph{really} know what you are doing, check the package
            code and look for the \texttt{geometry} class.
      \end{itemize}
\end{itemize}
\end{frame}

\section{Diagrams}

\begin{frame}{Pipelines with Icons}
\centering
\begin{tikzpicture}[
    icon/.style={inner sep=0pt},
    lbl/.style={align=center, font=\small, text width=2.2cm},
    flow/.style={-{Stealth[length=3mm]}, line width=1.2pt, subtitlegray},
    node distance=1.8cm]
  \node[icon] (data) {\includegraphics[width=1cm]{assets/icons/database}};
  \node[lbl, below=2mm of data] {Raw data};
  \node[icon, right=of data] (model) {\includegraphics[width=1cm]{assets/icons/robot}};
  \node[lbl, below=2mm of model] {Model};
  \node[icon, right=of model] (out) {\includegraphics[width=1cm]{assets/icons/filter}};
  \node[lbl, below=2mm of out] {Refined data};
  \draw[flow] (data) -- (model);
  \draw[flow] (model) -- (out);
\end{tikzpicture}

\vspace{\baselineskip}
\raggedright
Icons are small pictures placed in Ti\emph{k}Z nodes, connected by arrows.
They can be PDF/PNG files (e.g.\ from \hrefcol{https://www.flaticon.com}{Flaticon})
or glyphs of an icon font such as \texttt{fontawesome5}.
\end{frame}

\begin{frame}[fragile]{Pipelines with Icons}
\framesubtitle{The code}
\begin{block}{Basic pipeline}
\footnotesize
\begin{verbatim}
% preamble: \usetikzlibrary{positioning,arrows.meta}
\begin{tikzpicture}[node distance=1.8cm,
    flow/.style={-{Stealth}, line width=1.2pt, subtitlegray}]
  \node (data) {\includegraphics[width=1cm]{assets/icons/database}};
  \node[below=2mm of data] {Raw data};
  \node[right=of data] (model) {\includegraphics[width=1cm]{assets/icons/robot}};
  \node[below=2mm of model] {Model};
  \draw[flow] (data) -- (model);
\end{tikzpicture}
\end{verbatim}
\end{block}
With an icon font: \verb|\usepackage{fontawesome5}| and
\verb|\node[text=maincolor,font=\Huge] (data) {\faDatabase};|
\end{frame}

\begin{frame}{A More Elaborated Pipeline}
\framesubtitle{Revealed step by step}
\centering
\begin{tikzpicture}[
    icon/.style={inner sep=0pt},
    lbl/.style={align=center, font=\scriptsize, text width=1.6cm},
    flow/.style={-{Stealth[length=2.5mm]}, line width=1pt, subtitlegray},
    note/.style={font=\scriptsize, text=subtitlegray},
    node distance=1.25cm]
  % 1: raw data
  \node[icon] (raw) {\includegraphics[width=0.85cm]{assets/icons/database}};
  \node[lbl, below=1.5mm of raw] (rawl) {Raw data};
  % 2: automatic stage
  \node[icon, right=of raw] (clean) {{\visible<2->{\includegraphics[width=0.85cm]{assets/icons/broom}}}};
  \node[visible on=<2->, lbl, below=1.5mm of clean] (cleanl) {Cleaning};
  \node[icon, right=of clean] (model) {{\visible<2->{\includegraphics[width=0.85cm]{assets/icons/robot}}}};
  \node[visible on=<2->, lbl, below=1.5mm of model] (modell) {LLM judge};
  \node[icon, right=of model] (filt) {{\visible<2->{\includegraphics[width=0.85cm]{assets/icons/filter}}}};
  \node[visible on=<2->, lbl, below=1.5mm of filt] (filtl) {Filtering};
  \draw[visible on=<2->, flow] (raw) -- (clean);
  \draw[visible on=<2->, flow] (clean) -- (model);
  \draw[visible on=<2->, flow] (model) -- node[above, note] {scores} (filt);
  \begin{scope}[on background layer]
    \node[visible on=<2->, draw=maincolor, dashed, rounded corners, inner sep=6pt,
          fit=(clean) (cleanl) (filt) (filtl),
          label={[note, text=maincolor, visible on=<2->]above:Automatic stage}] (auto) {};
  \end{scope}
  % 3: human check, with a loop back to the model
  \node[icon, right=of filt] (human) {{\visible<3->{\includegraphics[width=0.85cm]{assets/icons/user}}}};
  \node[visible on=<3->, lbl, below=1.5mm of human] (humanl) {Human check};
  \draw[visible on=<3->, flow] (filt) -- (human);
  \draw[visible on=<3->, flow, dashed] (humanl.south) -- ++(0,-0.75cm)
    -| node[pos=0.25, below, note] {disagreements: refine the prompt} (modell.south);
  % 4: final dataset
  \node[icon, above=1.1cm of human] (final) {{\visible<4->{\includegraphics[width=0.85cm]{assets/icons/check}}}};
  \node[visible on=<4->, lbl, right=1.5mm of final, text width=1.4cm, align=left] {Final dataset};
  \draw[visible on=<4->, flow] (human) -- (final);
\end{tikzpicture}

\vspace{\baselineskip}
\raggedright\small
\only<1>{Parts are revealed at each click with overlay specifications (\texttt{visible on=<2->}, see the preamble).}%
\only<2>{\texttt{fit} draws a box around a group; \texttt{on background layer} keeps it behind the icons.}%
\only<3>{\texttt{|-} and \texttt{-|} draw right-angled paths, here for a feedback loop.}%
\only<4>{See the source of this slide in \texttt{beamersapienza.tex}.}
\end{frame}


\section{Summary}

\begin{frame}
\frametitle{Good Luck!}
\begin{itemize}
\item Enough for an introduction! You should know enough by now
\item If you have any suggestions or corrections, feel free to contribute to the \hrefcol{https://github.com/andrea-gasparini/sapienza-beamer-template}{GitHub repository}!
You can \hrefcol{https://github.com/andrea-gasparini/sapienza-beamer-template/issues/new/choose}{open an issue} or \hrefcol{https://github.com/andrea-gasparini/sapienza-beamer-template/fork}{fork the project} and directly propose your changes with a Pull Request. 
\end{itemize}
\end{frame}

\backmatter
\end{document}
